% 致读者
\chapter*{A Note to the Reader}
\addcontentsline{toc}{chapter}{A Note to the Reader}
\addtocontents{toc}{\vskip 1.2cm}
\vspace{6.5cm}
Two matters require comment---the exercises and the examples.

Working problems is a crucial part of learning mathematics. 
No one can learn topology merely by poring over the definitions, theorems, 
and examples that are worked out in the text.One must work part of it out for 
oneself. To provide that opportunity is the purpose of the exercises.

They vary in difficulty, with the easier ones usually given first. Some are 
routine verifications designed to test whether you have understood the definitions 
or examples of the preceding section.Others are less routine.
You may, for instance, be asked to generalize a theorem of the text. 
Although the result obtained may be interesting in its own right, 
the main purpose of such an exercise is to encourage you to work carefully 
through the proof in question, mastering its ideas thoroughly---more thoroughly 
(I hope!) than mere memorization would demand.

Some exercises are phrased in an "open-ended" fashion. Students often find 
this practice frustrating. When faced with an exercise that asks, 
``Is every regular Lindelöf space normal?" they respond in exasperation, 
``I don't know what I'm supposed to do! Am I suppose to prove it or find a 
counterexample or what?" But mathematics (outside textbooks) is usually like 
this. More often than not, all a mathematician has to work with is a conjecture 
or question, and he or she doesn't know what the correct answer is. 
You should have some experience with this situation.

A few exercises that are more difficult than the rest are marked with asterisks.
But none are so difficult but that the best student in my class can usually solve them.

Another important part of mastering any mathematical subject is acquiring a 
repertoire of useful examples. One should, of course, come to know those major 
examples from whose study the theory itself derives, and to which the important 
applications are made. But one should also have a few counterexamples at hand 
with which to test plausible conjectures.

Now it is all too easy in studying topology to spend too much time dealing with 
"weird counterexamples." Constructing them requires ingenuity and is often great 
fun. But they are not really what topology is about. Fortunately, one does not 
need too many such counterexamples for a first course; there is a fairly short list 
that will suffice for most purposes. Let me give it here:

{\advance\leftskip-1em
\hangindent=4em   % 悬挂缩进宽度（从第2行起左缩进的距离）
\hangafter=1      % 从第1行之后开始悬挂（即第1行正常，后续行缩进）
\({\mathbb{R}}^{J}\) the product of the real line with itself, in the product, uniform, and box topologies.

\({\mathbb{R}}_{\ell }\) the real line in the topology having the intervals \(\lbrack a,b)\) as a basis.

\({S}_{\Omega }\) the minimal uncountable well-ordered set.

\({I}_{o}^{2}\) the closed unit square in the dictionary order topology.
}
These are the examples you should master and remember; they will be exploited again and again.